% ============================================================
% §2  Related Work  (S version, ICLR 9-page) -- compressed from
% drafts/s2_related.tex (1075w -> <=430w). Writer, session (S2 compress).
% Drift contract (STORY S2 + R-rules):
%   * GradProm differentiation paragraph PRESERVED verbatim-welded
%     (MI lower bound + routing vs. gradient repair; R12).
%   * query-for-retake GAP kept (C2); IB block kept light; QC-VIB one footnote (R13).
%   * BMVC cited not reused (R10); method names anonymised (R4).
%   * R1/R2/R5 honoured: derive-style verbs, no overclaim, no leaderboard wording.
% All \citep keys carried over verbatim from the long draft.
% ============================================================

\section{Related Work}
\label{sec:related}

\textbf{Enhancement and the perception--diagnosis gap.}
Strong deterministic restoration backbones
(Real-ESRGAN~\citep{wang2021realesrgan}, Restormer~\citep{zamir2022restormer},
NAFNet~\citep{chen2022simple}, MIRNet-v2~\citep{zamir2022mirnetv2},
SwinIR~\citep{liang2021swinir}, Uformer~\citep{wang2022uformer}) optimise
distortion or perceptual fidelity, but high fidelity need not preserve
diagnosis. Generative super-resolution can even synthesise pigment-network or
vascular structure that was never present, a dermoscopy safety hazard; we
therefore use deterministic restoration and analyse this failure mode in
\S\ref{sec:discussion}. GradProm~\citep{gradprom2025} documents the same
gap in the skin-lesion (ISIC) domain and repairs gradients during training.
We share its diagnosis but differ in guarantee and scope:
\begin{quote}
Unlike GradProm, which corrects gradients during training, we provide (i) a
mutual-information lower bound guarantee for diagnosis preservation
(Lemma~\ref{lem:dp}), and (ii) a routing decision including query-for-retake.
GradProm treats the symptom by repairing gradients; we bound the information
loss and decide when enhancement is unsafe.
\end{quote}
We cite GradProm as prior evidence of the gap and differentiate on
\emph{mechanism}, running no head-to-head comparison and claiming no superiority.
Across these enhancement and restoration families, none \emph{preserves the
diagnostic signal under a guarantee}: gradient-level repair treats the symptom,
whereas we bound the information loss (Lemma~\ref{lem:dp}) and route on it.

\textbf{Perception--distortion trade-offs.}
The pull between perceptual quality and a downstream criterion is structural:
Blau and Michaeli~\citep{blau2018perception} bound the perception--distortion
trade-off, and Cohen et al.~\citep{cohen2018distribution} show that
distribution-matching translation can add or remove pathology-bearing features.
We instead make \emph{diagnostic} fidelity the explicitly preserved quantity,
with a mutual-information lower bound (Lemma~\ref{lem:dp}) characterising the
moderate-degradation window in which enhancement costs no diagnostic information.

\textbf{Selective prediction, deferral, and the recapture gap.}
Selective classification abstains below a confidence threshold
\citep{geifman2017selective,geifman2019selectivenet}, and learning-to-defer
hands the case to a clinician; both treat the input as fixed. Medical
image-quality work instead flags poor capture and prompts re-photography
(teledermatology~\citep{vodrahalli2021teledermatology,jalaboi2023explainable},
point-of-care ultrasound~\citep{zhao2020lungus}), supplying clinical
motivation but leaving the retake decision decoupled from any diagnostic-risk
estimate. Our query-for-retake channel closes this gap: the agent
(\S\ref{sec:agent}) routes on a quality estimate and, when the image is too
degraded for enhancement to be safe, requests re-acquisition --- a re-acquisition
loop that is, to our knowledge, the first to be bounded by a local diagnostic-risk
guarantee and coupled jointly with the diagnostic decision, rather than triggered by
a quality flag in isolation. We delimit this claim against adjacent lines of work that
might appear to subsume it: active perception and next-best-view planning
choose a sensing action to reduce task uncertainty, active sensing and active feature
acquisition select which measurements to obtain under a budget, and human-in-the-loop
recapture asks a clinician to re-photograph; none of these couples the re-acquisition
decision to a per-input diagnostic-risk estimate with a local risk bound, and so none
falls within the risk-bounded, diagnosis-coupled routing framework we study here. We
make no claim that this routing raises net benefit globally; \S\ref{sec:exp-dca}
reports where it does and does not help.

\textbf{Information bottleneck.}
The information bottleneck~\citep{tishby2000information} and its variational
form~\citep[VIB;][]{alemi2017deep} compress an input while retaining
task-relevant information, the lens for our diagnostic-loss analysis
(Lemma~\ref{lem:dp}).%
\footnote{One variant conditions the variational prior on a per-input quality
scalar; it yielded no diagnostic advantage over standard VIB, so we do not rely
on it---any off-the-shelf per-input uncertainty estimate suffices for routing,
and its derivation is retained in the appendix as framework only.}
